% https://confluence.cern.ch/display/IC/Absolute+energy+and+machine+length+calibration+in+the+SPS



\documentclass[a4paper,
               %boxit,        % check whether paper is inside correct margins
               %titlepage,    % separate title page
               %refpage       % separate references
               biblatex,     % biblatex is used
               keeplastbox,   % flushend option: not to un-indent last line in References
               %nospread,     % flushend option: do not fill with whitespace to balance columns
               hyphens,      % allow \url to hyphenate at "-" (hyphens)
               %xetex,        % use XeLaTeX to process the file
               %luatex,       % use LuaLaTeX to process the file
               ]{jacow}
%
% ONLY FOR \footnote in table/tabular
%
\usepackage{pdfpages,multirow,ragged2e} %
%
% CHANGE SEQUENCE OF GRAPHICS EXTENSION TO BE EMBEDDED
% ----------------------------------------------------
% test for XeTeX where the sequence is by default eps-> pdf, jpg, png, pdf, ...
%    and the JACoW template provides JACpic2v3.eps and JACpic2v3.jpg which
%    might generates errors, therefore PNG and JPG first
%
\makeatletter%
	\ifboolexpr{bool{xetex}}
	 {\renewcommand{\Gin@extensions}{.pdf,%
	                    .png,.jpg,.bmp,.pict,.tif,.psd,.mac,.sga,.tga,.gif,%
	                    .eps,.ps,%
	                    }}{}
\makeatother

% CHECK FOR XeTeX/LuaTeX BEFORE DEFINING AN INPUT ENCODING
% --------------------------------------------------------
%   utf8  is default for XeTeX/LuaTeX
%   utf8  in LaTeX only realises a small portion of codes
%
\ifboolexpr{bool{xetex} or bool{luatex}} % test for XeTeX/LuaTeX
 {}                                      % input encoding is utf8 by default
 {\usepackage[utf8]{inputenc}}           % switch to utf8

\usepackage[USenglish]{babel}

%
% if BibLaTeX is used
%
\ifboolexpr{bool{jacowbiblatex}}%
 {%
  \addbibresource{WEPM001_ref.bib}
 }{}
\listfiles

%%
%%   Lengths for the spaces in the title
%%   \setlength\titleblockstartskip{..}  %before title, default 3pt
%%   \setlength\titleblockmiddleskip{..} %between title + author, default 1em
%%   \setlength\titleblockendskip{..}    %afterauthor, default 1em

\begin{document}

\title{Spin tracking and equilibrium polarization computation in Xsuite: implementation and benchmarks}

\author{
G. Iadarola,
R. De Maria,
K. Hock,
J. Keintzel,
S. Lopaciuk,
J. Wenninger
}
\maketitle

%
\begin{abstract}
Precise modelling of spin dynamics is essential for the energy-calibration programme of the FCC-ee collider, where resonant depolarization measurements require accurate predictions of polarization buildup and of the spin-tune response to machine errors. We present spin modeling capabilities recently implemented in the Xsuite simulation framework, enabling high-performance six-dimensional tracking of particle trajectories together with their spin motion. The module includes the integration of the Thomas–BMT equation and a linearized computation of the invariant spin field (ISF), obtained from the one-turn linear map. Radiative spin effects are incorporated through the Sokolov–Ternov mechanism and its generalizations, allowing the evaluation of equilibrium polarization and of its buildup time.
The new tool has been benchmarked against the BMAD code for representative lattices, covering both spin-closed-orbit properties and long-term polarization evolution.
\end{abstract}


\end{document}







% \documentclass{article}
% \usepackage{graphicx} % Required for inserting images

% \title{}
% \author{giovanni.iadarola }
% \date{November 2025}

% \begin{document}

% \maketitle

% \section{Introduction}

% \end{document}



